\begingroup\fontsize{9.3}{10.7}\selectfont\setlength{\parskip}{1pt}
\section{Notation des réalisations gravitationnelles}
La notation distingue les objets de la construction, les opérateurs
de réponse et leurs réalisations physiques. Les produits intérieurs positifs
des extensions et la métrique lorentzienne ont des fonctions différentes.
\renewcommand{\arraystretch}{1.0}\begin{longtable}{@{}p{.23\linewidth}p{.72\linewidth}@{}}
\toprule
Symbole & Objet et fonction\\ \midrule
\endhead
\(G_a,\hbar_a\) & Sections réciproques de gravitation et d'action.
Leur produit détermine l'aire invariante
\(\ell_P^2=\hbar_aG_a/c^3=L_\alpha^2\) dans la réalisation radiale.\\
\(L_*,L_\alpha=L\) & Section longitudinale de référence et longueur construite \(L_\alpha=(5759/23040)\alpha^{16}L_*\).\\
\(E_L,t_P,t_0\) & Énergie de retour \(E_L=\hbar_{\rm ret}c/L_\alpha\), temps \(t_P=L_\alpha/c\) et pas \(t_0=\pi_{\rm HMT}t_P/54\).\\
\(R_X,\Lambda_X\) & Lecteur radial positif et longueur conjuguée du
même caractère ; \(R_X\Lambda_X=L_\alpha^2I\).\\
\(R_h,\zeta_h\) & Rayon aréolaire d'horizon et coordonnée \(R_h/L_\alpha=\zeta_h>0\). Ils conservent la condition géométrique de l'horizon.\\
\(\chi_{\rm rad}=L_\alpha/E_L\) & Proportion dimensionnelle entre rayon
et énergie. Ses unités et sa fonction diffèrent de celles du registre \(K\)
et de \(\kappa=8\pi G/c^4\).\\
\(\gamma_{\rm HMT}\) & Évaluation de la fonctionnelle hexade--octade,
réutilisée comme coefficient de l'opérateur de Holst.\\
\(k_B,\Theta\) & Coefficient de conversion informationnelle et température dans la section thermique déclarée.\\
\(P_5,\Gamma_5,\Lambda_n\) & Projection sur les cinq coordonnées,
isométrie tensorielle et réduction transversale de rang deux.\\
\(c_s,c_g\) & Cocycle de mémoire translationnelle et composante de retour du transport orienté.\\
\(\beta_m^\square,\Lambda_m^\square\) & Soudure d'écran et
réponse énergétique, avec son identité quadratique et son échelle.\\
\(\mathsf C,\mathsf L,f_{\min}\) & Incidence transportée,
gramien et extension d'énergie minimale pour la donnée de frontière.\\
\(s,\sigma,\mathsf M_{\rm cel}\) & Covecteur de la différentielle, courant matériel et accouplement qui convertit ses coordonnées.\\
\(e,\omega,\mathring\omega,K,T\) & Cotétrade, connexion, connexion de
Levi--Civita, contorsion et torsion.\\
\(\mathsf P_\gamma,\mathsf L_e\) & Opérateur interne de Holst et
application linéaire de la torsion à l'équation de connexion.\\
\(Q,F,\Delta S\) & Terme gravitationnel quadratique, action matérielle
et correction effective après élimination de la contorsion.\\
\(\mathsf H\) & Hessienne de l'action composée qui contrôle la continuation locale d'une branche torsionnelle.\\
\(M_n,a,a_{\rm ref}\) & Mémoire, facteur d'échelle et normalisation
de référence de sa réalisation homogène.\\
\(s_0,\rho_0,w\) & Amplitude torsionnelle, densité matérielle de
référence et paramètre barotrope du domaine de rebond.\\
\(\varrho,V_c,\widehat{\mathscr Q}_{\rm sp}\) & État matériel normalisé, volume comobile et observable quartique normalement ordonné ; ils déterminent la fonctionnelle \(s_0[\varrho,V_c]\).\\
\(C_\gamma\) & Coefficient de l'interaction effective spinorielle
obtenu par élimination de la contorsion.\\
\(\tau,t_{\rm b}\) & Durée caractéristique et instant du minimum
de la solution homogène de poussière.\\
\(F_0,R,a_{\rm b}\) & Fonction de Friedmann, perturbation et seuil transversal de l'évolution homogène.\\
\(T_\Lambda,T_0\) & Parties de trace et sans trace du tenseur résiduel ;
leurs divergences enregistrent l'échange entre les secteurs.\\
\(\mathsf{G}_{\rm ret},\mathsf{G}_{\rm adv}\) & Propagateurs retardé et avancé.
Ici, la lettre \(G\) désigne des opérateurs, et non la constante gravitationnelle.\\
\(P_{\rm abs}\) & Projection sur les sources qui satisfont l'égalité des deux lectures causales à la frontière.\\
\bottomrule
\end{longtable}

\par\endgroup
